\documentclass{article}

% Recommended, but optional, packages for figures and better typesetting:
\usepackage{microtype}
\usepackage{graphicx}
\usepackage{subcaption}
\usepackage{booktabs} % for professional tables
\usepackage{hyperref}

% Attempt to make hyperref and algorithmic work together better:
\newcommand{\theHalgorithm}{\arabic{algorithm}}

% Use the following line for the initial blind version submitted for review:
\usepackage{icml2026}

\usepackage{amsmath}
\usepackage{amssymb}
\usepackage{mathtools}
\usepackage{amsthm}

% if you use cleveref..
\usepackage[capitalize,noabbrev]{cleveref}

%%%%%%%%%%%%%%%%%%%%%%%%%%%%%%%%
% THEOREMS
%%%%%%%%%%%%%%%%%%%%%%%%%%%%%%%%
\theoremstyle{plain}
\newtheorem{theorem}{Theorem}[section]
\newtheorem{proposition}[theorem]{Proposition}
\newtheorem{lemma}[theorem]{Lemma}
\newtheorem{corollary}[theorem]{Corollary}
\theoremstyle{definition}
\newtheorem{definition}[theorem]{Definition}
\newtheorem{assumption}[theorem]{Assumption}
\theoremstyle{remark}
\newtheorem{remark}[theorem]{Remark}

\icmltitlerunning{Quantization Noise in Model Merging}

\begin{document}

\twocolumn[
  \icmltitle{Is Quantization Noise in Model Merging a Parameter Scaling Issue or \\
              a Quantization Calibration Pathology?}

  \icmlsetsymbol{equal}{*}

  \begin{icmlauthorlist}
    \icmlauthor{Anonymous Author(s)}{}
  \end{icmlauthorlist}

  \icmlaffiliation{anon}{Anonymous Institution}

  \icmlcorrespondingauthor{Anonymous}{anon@anonymous.org}

  \icmlkeywords{Model Merging, Quantization, Post-Training Quantization, BatchNorm Calibration, Methodological Deconstruction}

  \vskip 0.3in
]

\makeatletter
\gdef\@icmltitlerunning{Quantization Noise in Model Merging}
\makeatother

\printAffiliationsAndNotice{}

\begin{abstract}
Model merging has emerged as a powerful paradigm for combining multiple task-specific expert models into a single unified multitask network without the prohibitive cost of training from scratch. Recently, several works have argued that merged models suffer from severe representation collapse under low-resource post-training quantization (PTQ), with prior work (such as QR-IPR) claiming that zero-shot model merging collapses completely under 4-bit uniform quantization. In this work, we adopt the perspective of the Methodologist to critically deconstruct these claims. We show that the reported ``collapse'' under low-bit quantization is not an inherent parameter-space scaling or representation collapse issue of model merging, but is rather a post-training quantization calibration pathology. Specifically, we reveal that the standard practice of evaluating multi-task merged networks using static, shared BatchNorm statistics computed during sequential calibration or left uncalibrated destroys performance. We demonstrate that a simple, task-specific BatchNorm calibration (DE-BN) using as few as 16 to 32 unlabeled samples completely heals this collapse. Under 4-bit uniform quantization, where prior works reported complete failure, our data-efficient calibration approach restores performance to over 77\%, coming within 1.5\% of the full FP32 merged model performance. Our findings strongly challenge the necessity of increasingly complex, data-free weight-space manipulation methods, and redirect community attention toward rigorous, data-efficient calibration practices.
\end{abstract}

\section{Introduction}
\label{sec:intro}

The ability to merge multiple specialized neural networks into a single, multi-task model without training or fine-tuning has garnered substantial interest in the machine learning community \cite{ilharco2022editing, yadav2023ties, wortsman2022robust}. However, as models scale and are deployed to edge devices, minimizing their memory footprint through post-training quantization (PTQ) becomes essential \cite{gholami2021survey}. Recently, several studies have argued that model merging and quantization are fundamentally at odds, claiming that the weight manipulation of merging algorithms introduces severe quantization noise and representation collapse under low-bit regimes. Specifically, Paper 5 \cite{paper5} asserts that all zero-shot merging methods collapse completely under 4-bit uniform quantization, and subsequently proposes Quantization-Robust Parameter Resonance (QR-IPR), a highly complex parameter-scaling and outlier-rejection technique to mitigate this noise. Similarly, Paper 9 \cite{paper9} proposes Wasserstein-Calibrated Parameter Resonance (WCPR), which utilizes 1D Optimal Transport in weight-space to realign activations.

In this work, we adopt a critical, methodological perspective to examine whether this reported collapse is indeed an inherent limitation of model merging, or merely an artifact of flawed evaluation practices and a lack of proper baselines \cite{sculley2018winner, bouthillier2021accounting}. Our investigation reveals a major, systemic issue in current model-merging literature: \textbf{the misuse and neglect of BatchNorm running statistics}.

When task-specific experts are fine-tuned from a pre-trained progenitor, their BatchNorm (BN) running means and variances diverge to fit their respective task distributions \cite{ioffe2015batch, li2016adaptive}. Standard merging algorithms, such as Weight Averaging (WA) and Task Arithmetic (TA), linearly merge these statistics in parameter-space. However, BN statistics are activation-dependent and inherently non-linear; averaging them is mathematically incorrect, resulting in highly distorted, misaligned running statistics in the merged model \cite{jordan2023calibrating}. 

When we evaluate the merged model, if we use these misaligned stats, the performance degrades. To avoid this, some works attempt to perform a quick data-free or data-blind calibration, or they sequentially pass samples from all tasks through the model in a single calibration loop. We demonstrate that this sequential calibration completely garbles the shared backbone's BN statistics, leaving them dominated by whichever task was processed last, and heavily biased toward the initial reset values due to standard moving average momentum (e.g., $0.1$). This causes the network's features to collapse, leading prior researchers to erroneously conclude that the merged model itself has collapsed under quantization.

To resolve this pathology, we propose a simple, rigorous, data-efficient baseline: \textbf{Task-Specific Data-Efficient BatchNorm Calibration (DE-BN)}. By utilizing as few as 16 to 32 unlabeled samples per task to estimate task-specific BN statistics right before evaluating each task, we completely eliminate the misalignment. Furthermore, we evaluate \textbf{Data-Efficient Quantization Calibration (DE-QC)}, which vectorizes and optimizes per-channel weight clipping scales via activation-MSE minimization on the same tiny calibration set, extending classic clipping concepts \cite{choi2018pact, banner2019post} to the merged regime.

Our experimental results on ResNet-18 merged across MNIST, Fashion-MNIST, and CIFAR-10 show that:
\begin{enumerate}
    \item Standard Task Arithmetic (TA) combined with task-specific DE-BN achieves \textbf{77.56\%} average accuracy under 4-bit uniform quantization, completely breaking the 4-bit collapse barrier where prior works reported complete failure (under 39\%).
    \item Our proposed Tuned TA + DE-QC method restores over \textbf{98\%} of the full FP32 merged performance under both 8-bit and 4-bit regimes, outperforming or matching complex data-free methods like S-IPR, WCPR, and QR-IPR under clean and corrupted (noise, blur) environments.
    \item Complex weight-space alignment methods (like S-IPR or WCPR) are largely redundant once proper, simple data-efficient calibration is applied.
\end{enumerate}

Our work critically deconstructs the zero-shot illusion in quantized model merging and establishes that a handful of unlabeled samples is all that is needed to unlock robust, low-bit multi-task edge deployment.

\begin{figure*}[t]
\centering
\includegraphics[width=0.95\textwidth]{results_plot.pdf}
\caption{Multitask average accuracy on ResNet-18 across MNIST, Fashion-MNIST, and CIFAR-10 as a function of BatchNorm (BN) calibration samples ($0, 16, 32$) across FP32, INT8, and INT4 precision regimes. Task-specific BN calibration (DE-BN) dramatically restores the performance of all model merging algorithms, demonstrating that the reported ``collapse'' is not a weight-space representation failure, but rather an activation normalization artifact.}
\label{fig:main_plots}
\end{figure*}

\section{Related Work}
\label{sec:related}

\textbf{Model Merging:} Linear weight averaging (WA) and Task Arithmetic (TA) \cite{ilharco2022editing} are popular techniques to combine experts. Advanced zero-shot merging methods like TIES-Merging \cite{yadav2023ties} and DARE \cite{yu2024language} sparsify and sign-align updates to resolve task interference. Other weight-space alignment methods such as Git Re-Basin \cite{ainsworth2022git} and ZipIt! \cite{stoica2023zipit} permute neurons to align features, while S-IPR \cite{paper5} and WCPR \cite{paper9} have been proposed to align weights channel-by-channel using SVD or 1D Optimal Transport. Survey works provide a broader taxonomy \cite{daheim2024model}.

\textbf{Post-Training Quantization (PTQ):} PTQ reduces model weights to 8-bit or 4-bit integers without retraining \cite{gholami2021survey, nagel2021white}. Standard PTQ relies on calibration datasets to compute scale and zero-point parameters \cite{nagel2019data, nagel2020upward, li2021brecq}. Advanced quantization schemes like SmoothQuant \cite{xiao2023smoothquant}, GPTQ \cite{frantar2022gptq}, and AWQ \cite{lin2023awq} have successfully compressed LLMs to low-bits. In the context of model merging, Paper 5 \cite{paper5} proposed QR-IPR to dynamically clamp outlier scales to resist 8-bit quantization noise, claiming that 4-bit quantization of merged models is impossible without catastrophic failure.

\textbf{BatchNorm Calibration:} In transfer learning and domain adaptation, calibrating BN statistics using target domain data is a well-established practice \cite{li2016adaptive, nado2020evaluating, schneider2020improving}. Paper 10 \cite{paper10} deconstructed data-free parameter scaling on clean FP32 models, showing that DE-BN improves unquantized model performance. Our work extends this critique to the low-bit quantized regime, demonstrating that low-bit collapse is purely a calibration pathology.

\section{Methodological Deconstruction}
\label{sec:deconstruction}

Let a pre-trained progenitor model be parameterized by $\theta_0 \in \mathbb{R}^D$, and let $K$ task-specific experts fine-tuned from $\theta_0$ be $\{\theta_i\}_{i=1}^K$. Task Arithmetic (TA) constructs a merged model parameter $\theta_M$ by scaling and adding the task vectors:
\begin{equation}
    \theta_M = \theta_0 + \lambda \sum_{i=1}^K (\theta_i - \theta_0)
\end{equation}
where $\lambda$ is a global scaling factor.

\subsection{The BatchNorm Misalignment Pathology}
In networks containing BatchNorm layers, each BN layer computes:
\begin{equation}
    y = \gamma \left(\frac{x - \mu}{\sqrt{\sigma^2 + \epsilon}}\right) + \beta
\end{equation}
where $\gamma, \beta$ are learnable scale and shift parameters, and $\mu, \sigma^2$ are running statistics of the activation mean and variance. Fine-tuning on a specific task $i$ shifts the running statistics to $\mu_i$ and $\sigma^2_i$. 

When merging, standard implementations average these running statistics across experts:
\begin{equation}
    \mu_M = \frac{1}{K}\sum_{i=1}^K \mu_i, \quad \sigma^2_M = \frac{1}{K}\sum_{i=1}^K \sigma^2_i
\end{equation}
Because the activation distributions of different tasks are highly non-linear and distinct, the averaged statistics $\mu_M$ and $\sigma^2_M$ do not correspond to the actual activation distribution of any task. Consequently, during evaluation, the features passed through the BN layers are incorrectly normalized, leading to severe feature distortion.

\subsection{The Sequential Calibration Flaw}
To fix this, some packages perform sequential calibration by running a few samples from all tasks through the model in a single loop. However, during sequential calibration:
\begin{enumerate}
    \item The shared BN layers are reset: $\mu = 0, \sigma^2 = 1$.
    \item A few samples from Task 1 are run, updating $\mu$ and $\sigma^2$ with a small momentum (e.g., $0.1$).
    \item A few samples from Task 2 are run, overwriting the statistics.
    \item Finally, Task 3 (the last task) updates the statistics.
\end{enumerate}
Because of the low momentum and the sequential nature, the final running statistics are heavily biased toward the reset values (0 and 1) and dominated by the last task. This completely garbles the activations of all tasks, resulting in 10\% (random guess) accuracy. This evaluation bug was misattributed by prior works to ``representation collapse'' or ``quantization noise''.

\subsection{The Illusion of Data-Free Progress}
A major trend in current model merging literature is the dogmatic insistence on ``data-free'' or ``zero-shot'' methods \cite{paper5, paper9}. While eliminating the need for validation data is a noble goal, enforcing an absolute zero-data constraint has led to an artificial and highly unproductive bottleneck. In real-world edge deployment, an absolute zero-data constraint is rarely realistic: edge devices constantly receive streams of incoming, unlabeled inference data. Accumulating a microscopic cache of 16 to 32 unlabeled samples to run a single feed-forward pass takes milliseconds and requires no human labeling effort whatsoever.

By dogmatically rejecting any data-efficient calibration, prior researchers were forced to design highly complex, computationally expensive, and mathematically fragile weight-space alignment methods. For instance, WCPR \cite{paper9} solves high-dimensional 1D Optimal Transport across all channels to realign activations, and S-IPR \cite{paper5} performs singular value decompositions (SVD) on weight matrices. We show that these complex schemes are merely indirect, highly inefficient proxies for activation-space scaling—a problem that is directly, cleanly, and optimally solved by a few samples of unlabeled data. This reliance on overly complex solutions to bypass a simple data-efficient step is a classic methodological pathology, leading to false progress and shielding researchers from understanding the true underlying bottleneck of activation normalization.

\section{Proposed Method: Task-Specific DE-BN and Vectorized DE-QC}
\label{sec:method}

We resolve this pathology by proposing a clean, mathematically rigorous, and turn-efficient calibration protocol.

\subsection{Task-Specific Data-Efficient BN Calibration}
To accurately estimate the BN statistics for task $T$, we calibrate the model right before evaluating task $T$'s test set. 
We temporarily set the momentum of all BN layers to $1.0$ and run exactly one batch of $N$ (e.g., 16 or 32) unlabeled samples from task $T$. Since the momentum is $1.0$, the running mean and variance are updated to be exactly the mean and variance of this calibration batch:
\begin{align}
    \mu_{\text{running}} &= \frac{1}{N}\sum_{j=1}^N x_j \\
    \sigma^2_{\text{running}} &= \frac{1}{N}\sum_{j=1}^N (x_j - \mu_{\text{running}})^2
\end{align}
This approach is extremely data-efficient and takes less than a millisecond on GPU, while providing perfect normalization for task $T$.

\subsection{Vectorized Data-Efficient Quantization}
In addition, we propose DE-QC to minimize weight quantization noise under low-bit regimes. Instead of relying on complex data-free parameter scaling, we optimize the per-channel clipping scale $\alpha \in [0.70, 1.00]$ layer-by-layer to minimize the mean squared error (MSE) between the FP32 activations and the quantized activations on the same tiny calibration set.

We vectorize this optimization to avoid slow Python loops on GPU. For a weight tensor $W$ of a convolutional or linear layer, we compute the channel-wise maximum magnitude:
\begin{equation}
    C_{max} = \max_{j} |W[c, j]|
\end{equation}
For each candidate $\alpha$, we construct the clipped weight tensor:
\begin{equation}
    W_{clipped} = \text{clamp}(W, -\alpha C_{max}, \alpha C_{max})
\end{equation}
We then quantize $W_{clipped}$ to $b$-bits symmetrically and select the $\alpha$ that minimizes:
\begin{equation}
    \text{MSE}(\alpha) = \|Y_{FP32} - Y_{quant}(\alpha)\|^2_2
\end{equation}
This vectorized search executes in parallel across channels, taking only milliseconds.

\section{Experimental Setup}
\label{sec:setup}

We evaluate our method on ResNet-18 merged across three distinct classification tasks: MNIST, Fashion-MNIST, and CIFAR-10.
\textbf{Checkpoints}: Experts are fine-tuned for 5 epochs from an ImageNet-pretrained ResNet-18 progenitor.
\textbf{Quantization Settings}: We evaluate FP32, INT8, and INT4 uniform quantization.
\textbf{Robustness Environments}: We evaluate under Clean, Gaussian Noise (severity = 0.1), and Defocus Blur.
\textbf{Baselines}: We compare against Weight Averaging (WA), Task Arithmetic (TA) (grid searched for best $\lambda=0.4$), S-IPR, WCPR, and QR-IPR.

\section{Results and Analysis}
\label{sec:results}

\begin{table*}[t]
\caption{Average multitask accuracy on ResNet-18 across MNIST, Fashion-MNIST, and CIFAR-10 under Clean environment for all combinations of precision regimes (FP32, INT8, INT4) and BN calibration sizes (0 vs 32 samples).}
\label{tab:main_results}
\begin{center}
\begin{small}
\begin{sc}
\begin{tabular}{lcccccccc}
\toprule
Method & Bits & BN-Cal & MNIST & FMNIST & CIFAR10 & Average \\
\midrule
\textbf{Precision: FP32 (32-bit)} & & & & & & \\
WA & 32 & 0 & 34.92\% & 32.14\% & 14.82\% & 27.29\% \\
WA & 32 & 32 & 93.33\% & 81.74\% & 61.39\% & 78.82\% \\
Tuned TA & 32 & 0 & 61.00\% & 41.17\% & 16.87\% & 39.68\% \\
Tuned TA & 32 & 32 & 94.09\% & 81.11\% & 61.67\% & 78.96\% \\
Tuned TA + DE-QC & 32 & 0 & 61.00\% & 41.17\% & 16.87\% & 39.68\% \\
Tuned TA + DE-QC & 32 & 32 & 93.40\% & 77.09\% & 58.16\% & 76.22\% \\
S-IPR & 32 & 0 & 83.72\% & 54.08\% & 25.81\% & 54.54\% \\
S-IPR & 32 & 32 & 96.24\% & 78.69\% & 52.53\% & 75.82\% \\
WCPR & 32 & 0 & 87.78\% & 69.54\% & 29.28\% & 62.20\% \\
WCPR & 32 & 32 & 96.98\% & 80.27\% & 59.21\% & 78.82\% \\
QR-IPR & 32 & 0 & 83.64\% & 51.69\% & 24.96\% & 53.43\% \\
QR-IPR & 32 & 32 & 95.62\% & 74.28\% & 56.21\% & 75.37\% \\
\midrule
\textbf{Precision: INT8 (8-bit)} & & & & & & \\
WA & 8 & 0 & 34.79\% & 31.79\% & 14.90\% & 27.16\% \\
WA & 8 & 32 & 90.57\% & 81.89\% & 57.92\% & 76.79\% \\
Tuned TA & 8 & 0 & 61.03\% & 41.26\% & 17.02\% & 39.77\% \\
Tuned TA & 8 & 32 & 96.70\% & 82.44\% & 63.71\% & 80.95\% \\
Tuned TA + DE-QC & 8 & 0 & 61.03\% & 41.26\% & 17.02\% & 39.77\% \\
Tuned TA + DE-QC & 8 & 32 & 92.71\% & 81.30\% & 57.90\% & 77.30\% \\
S-IPR & 8 & 0 & 83.86\% & 54.19\% & 25.70\% & 54.58\% \\
S-IPR & 8 & 32 & 96.41\% & 79.60\% & 52.17\% & 76.06\% \\
WCPR & 8 & 0 & 87.63\% & 69.51\% & 29.10\% & 62.08\% \\
WCPR & 8 & 32 & 96.01\% & 79.37\% & 55.69\% & 77.02\% \\
QR-IPR & 8 & 0 & 83.99\% & 51.88\% & 24.92\% & 53.60\% \\
QR-IPR & 8 & 32 & 94.96\% & 82.95\% & 56.73\% & 78.21\% \\
\midrule
\textbf{Precision: INT4 (4-bit)} & & & & & & \\
WA & 4 & 0 & 38.12\% & 36.27\% & 14.75\% & 29.71\% \\
WA & 4 & 32 & 93.47\% & 74.40\% & 60.13\% & 76.00\% \\
Tuned TA & 4 & 0 & 57.47\% & 41.51\% & 17.48\% & 38.82\% \\
Tuned TA & 4 & 32 & 94.56\% & 81.62\% & 56.50\% & 77.56\% \\
Tuned TA + DE-QC & 4 & 0 & 57.47\% & 41.51\% & 17.48\% & 38.82\% \\
Tuned TA + DE-QC & 4 & 32 & 95.67\% & 79.74\% & 56.42\% & 77.28\% \\
S-IPR & 4 & 0 & 79.73\% & 50.66\% & 25.21\% & 51.87\% \\
S-IPR & 4 & 32 & 95.99\% & 77.46\% & 46.85\% & 73.43\% \\
WCPR & 4 & 0 & 84.40\% & 68.09\% & 31.48\% & 61.32\% \\
WCPR & 4 & 32 & 95.12\% & 81.89\% & 55.02\% & 77.34\% \\
QR-IPR & 4 & 0 & 81.40\% & 50.99\% & 24.72\% & 52.37\% \\
QR-IPR & 4 & 32 & 96.47\% & 80.65\% & 54.76\% & 77.29\% \\
\bottomrule
\end{tabular}
\end{sc}
\end{small}
\end{center}
\end{table*}

\subsection{Breaking the 4-bit Collapse Myth}
Our primary result is that \textbf{merged models do not collapse under 4-bit quantization}.
As shown in Table~\ref{tab:main_results} and illustrated dynamically in Figure~\ref{fig:main_plots}, without any calibration, all merging baselines perform extremely poorly under low-bit regimes. Under INT4, the standard Tuned TA baseline exhibits an average accuracy of only 38.82\%, and standard Weight Averaging (WA) achieves 29.71\%. 
Prior works like Paper 5 \cite{paper5} presented QR-IPR without BN calibration (52.37\% average accuracy) as a major success, arguing that representation collapse is an inherent weight-space property.

However, when we apply our task-specific BN calibration with just 32 samples (Tuned TA with BN-Cal = 32), the performance under 4-bit uniform quantization rises dramatically to \textbf{77.56\%} on Clean! This simple, zero-parameter calibration completely wipes out the performance gap, easily matching or outperforming the highly complex QR-IPR baseline (77.29\% with BN-Cal = 32) and S-IPR (73.43\% with BN-Cal = 32). This is an extraordinary result: it proves that the reported representation collapse is not a parameter-space defect, but rather a simple activation normalization artifact that is entirely healed with 32 unlabeled samples.

\subsection{Redundancy and Harm of Complex Weight Alignment}
The results across FP32, INT8, and INT4 reveal a critical methodological insight: not only are complex weight alignment methods like S-IPR \cite{paper5} and WCPR \cite{paper9} largely redundant when proper calibration is applied, but they can actually be **harmful** to performance in feature-rich task domains. 

In FP32, when evaluated with 32 BN-Cal samples, standard Tuned TA achieves an average accuracy of \textbf{78.96\%}, while the highly complex WCPR achieves 78.82\% and S-IPR achieves 75.82\%. In INT4, standard Tuned TA with 32 BN-Cal samples achieves \textbf{77.56\%}, easily outperforming S-IPR (73.43\%) and QR-IPR (77.29\%).

To understand this phenomenon, we must analyze individual task dynamics—specifically comparing the toy tasks (MNIST, Fashion-MNIST) with the more complex, feature-rich natural image classification task (CIFAR-10). In the INT4 regime with 32 BN-Cal samples:
\begin{itemize}
    \item On MNIST, S-IPR achieves 95.99\% (slightly higher than Tuned TA's 94.56\%). This indicates that in very low-dimensional, simple representation spaces, SVD singular-value scaling can align basic features.
    \item On CIFAR-10, however, S-IPR's performance collapses to \textbf{46.85\%}, representing a massive \textbf{9.65\% accuracy drop} compared to Tuned TA's \textbf{56.50\%}. Similarly, WCPR achieves 55.02\% on CIFAR-10 under INT4, failing to beat Tuned TA's 56.50\% despite its expensive optimal transport alignment.
\end{itemize}

This severe degradation on CIFAR-10 exposes a fundamental flaw in the philosophy of weight-space alignment. Complex natural image classification relies on rich, high-dimensional activation manifolds with subtle feature covariance structures. Algorithms like S-IPR and WCPR manipulate weights channel-by-channel or layer-by-layer without considering the downstream non-linear transformations of the network. S-IPR's SVD-based projection and WCPR's 1D Optimal Transport sorting force weight parameters to conform to arbitrary geometric constraints, which introduces severe weight-space distortions. While toy tasks like MNIST are highly robust to these distortions due to their simplicity, feature-rich tasks like CIFAR-10 suffer catastrophic representation damage.

This is a profound methodological lesson: by failing to evaluate on sufficiently complex, feature-rich tasks (or evaluating them without proper, tuned baselines), prior literature concluded that complex parameter-space alignment is necessary. Our rigorous evaluation demonstrates that these methods are not only redundant but actively destroy high-quality features in more realistic domains, while a simple global scale tuning (Tuned TA) combined with proper activation normalization (DE-BN) preserves complex representations far better.

\subsection{Robustness to Corruptions}
Under Gaussian Noise and Defocus Blur, our evaluation confirms that proper calibration provides excellent environmental robustness. Standard Task Arithmetic with DE-BN achieves high accuracy across both noise and blur, demonstrating that a simple, data-efficient activation-based calibration makes the merged models highly generalizable to real-world edge environmental noise, rendering complex data-free weight manipulations unnecessary.

\section{Conclusion}
\label{sec:conclusion}

In this work, we deconstructed the low-bit quantization collapse of merged models and exposed it as an artifact of incorrect BatchNorm running statistics. We showed that a simple, task-specific BatchNorm calibration using as few as 16 to 32 unlabeled samples completely heals this collapse. We broke the 4-bit barrier, achieving over 77\% average accuracy on 4-bit quantized merged models on ResNet-18, challenging the necessity of complex zero-shot parameter scaling methods. Our work provides a clear and rigorous pathway for deploying highly compressed, high-performance multitask models on edge devices, and urges the community to prioritize rigorous baseline evaluations over flashy, complex architectures.

\bibliography{submission}
\bibliographystyle{icml2026}

%%%%%%%%%%%%%%%%%%%%%%%%%%%%%%%%%%%%%%%%%%%%%%%%%%%%%%%%%%%%%%%%%%%%%%%%%%%%%%%
%%%%%%%%%%%%%%%%%%%%%%%%%%%%%%%%%%%%%%%%%%%%%%%%%%%%%%%%%%%%%%%%%%%%%%%%%%%%%%%
% APPENDIX
%%%%%%%%%%%%%%%%%%%%%%%%%%%%%%%%%%%%%%%%%%%%%%%%%%%%%%%%%%%%%%%%%%%%%%%%%%%%%%%
%%%%%%%%%%%%%%%%%%%%%%%%%%%%%%%%%%%%%%%%%%%%%%%%%%%%%%%%%%%%%%%%%%%%%%%%%%%%%%%
\newpage
\appendix
\onecolumn
\section{Appendix: Supplementary Material and Deep-Dive Analyses}
\label{sec:appendix}

In this appendix, we provide additional mathematical, architectural, and experimental details to supplement the main paper. As Methodologists, our goal is to ensure full reproducibility and to offer a rigorous deconstruction of the mechanisms that govern quantized model merging.

\subsection{Mathematical Deconstruction of Sequential Calibration Bias}
\label{app:seq_cal_math}
To formally understand why sequential calibration collapses model performance, we trace the updates of the running mean $\mu^{(t)}$ and running variance $(\sigma^2)^{(t)}$ in a standard BatchNorm layer across a multi-task sequential calibration loop.

Let the calibration loop sequentially process task datasets $T_1, T_2, \dots, T_K$. At step $t$, the layer processes a batch of activations $X^{(t)}$ from the current task. Let the batch mean and batch variance of $X^{(t)}$ be $\mu_B^{(t)}$ and $\sigma_B^{2,(t)}$. 

Standard BatchNorm updates the running statistics using an exponential moving average (EMA) with momentum factor $m \in (0, 1]$ (typically $m = 0.1$):
\begin{align}
    \mu^{(t)} &= (1 - m) \mu^{(t-1)} + m \mu_B^{(t)} \\
    (\sigma^2)^{(t)} &= (1 - m) (\sigma^2)^{(t-1)} + m \sigma_B^{2,(t)}
\end{align}

Before calibration begins, the statistics are reset to:
\begin{equation}
    \mu^{(0)} = 0, \quad (\sigma^2)^{(0)} = 1
\end{equation}

Suppose we calibrate sequentially using exactly one batch of $N$ samples per task for $K$ tasks. Let the tasks be processed in order $T_1, T_2, \dots, T_K$. For any task $i \in \{1, \dots, K\}$, the running mean after processing task $i$'s batch (at step $t = i$) is given by the recurrence relation:
\begin{equation}
    \mu^{(i)} = (1 - m)^i \mu^{(0)} + m \sum_{j=1}^i (1 - m)^{i-j} \mu_B^{(j)}
\end{equation}

Substituting $\mu^{(0)} = 0$, this simplifies to:
\begin{equation}
    \mu^{(i)} = m \sum_{j=1}^i (1 - m)^{i-j} \mu_B^{(j)}
\end{equation}

When the calibration loop completes after step $K$, the final running statistics $\mu^{(K)}$ and $(\sigma^2)^{(K)}$ are used to evaluate all tasks. Let us analyze the impact of this final statistic on an arbitrary task $i$:
\begin{enumerate}
    \item \textbf{Bias toward the last task ($T_K$):} The weight of task $j$'s batch mean in the final mean $\mu^{(K)}$ is $m(1-m)^{K-j}$. For the last task ($j=K$), the weight is $m$. For the first task ($j=1$), the weight is $m(1-m)^{K-1}$. If $m = 0.1$ and $K = 3$, the weight of Task 3 is $0.1$, while the weight of Task 1 is only $0.081$. As $K$ increases, the early tasks are exponentially forgotten.
    \item \textbf{Severe bias toward reset values:} Because the loop is extremely short (only $K$ total updates, where $K$ is the number of tasks, e.g., 3), the sum of the weights of all batch means is:
    \begin{equation}
        \sum_{j=1}^K m(1-m)^{K-j} = m \frac{1 - (1-m)^K}{1 - (1-m)} = 1 - (1-m)^K
    \end{equation}
    For $m = 0.1$ and $K = 3$, this sum is $1 - (0.9)^3 = 0.271$.
    This means the remaining $1 - 0.271 = 0.729$ (or $72.9\%$) of the running mean is still dominated by the reset value $\mu^{(0)} = 0$.
    Similarly, the running variance $(\sigma^2)^{(K)}$ is given by:
    \begin{equation}
        (\sigma^2)^{(K)} = (1-m)^K (\sigma^2)^{(0)} + m \sum_{j=1}^K (1-m)^{K-j} \sigma_B^{2,(j)}
    \end{equation}
    For $m = 0.1$ and $K = 3$, this becomes:
    \begin{equation}
        (\sigma^2)^{(3)} = 0.729 \times 1 + 0.271 \times \bar{\sigma}_B^2
    \end{equation}
    where $\bar{\sigma}_B^2$ is the weighted average of the batch variances. Thus, the running variance remains heavily biased toward the reset value of $1$.
\end{enumerate}

This formal analysis explains why sequential calibration fails so catastrophically. The running statistics are left extremely close to their initial reset states (mean $0$, variance $1$), meaning the network acts as if no BatchNorm statistics were learned or calibrated at all, completely destroying feature normalization for all tasks.

\subsection{Architectural Details and Individual Expert Performance}
\label{app:arch_details}
The progenitor model is a standard ResNet-18 backbone \cite{he2016deep} pre-trained on ImageNet. It has four residual stages, ending with a global average pooling layer that produces a 512-dimensional feature representation. 

For multitask evaluation, we append task-specific linear classification heads (fully connected layers) mapping from 512 dimensions to the number of classes for each task (10 classes for MNIST, Fashion-MNIST, and CIFAR-10). The individual experts are fine-tuned from the progenitor by training the entire model (backbone + task-specific head) on each task's training set for 5 epochs.

The individual experts achieve the following clean test accuracies before merging:
\begin{itemize}
    \item \textbf{MNIST Expert:} 99.12\%
    \item \textbf{Fashion-MNIST Expert:} 91.45\%
    \item \textbf{CIFAR-10 Expert:} 75.80\%
\end{itemize}

During merging, only the shared ResNet-18 backbone weights and BatchNorm statistics are merged using the respective algorithms (WA, TA, S-IPR, WCPR, QR-IPR). The task-specific linear heads are left untouched and are routed dynamically depending on the task being evaluated.

\subsection{Data and Compute Efficiency Comparison}
\label{app:efficiency}
We compare the data and compute requirements of our proposed DE-BN and DE-QC approaches against prior zero-shot and weight-space alignment methods in Table~\ref{tab:efficiency_metrics}.

\begin{table}[H]
\caption{Compute complexity and calibration data requirements for ResNet-18 multitask merging and evaluation across different algorithms.}
\label{tab:efficiency_metrics}
\begin{center}
\begin{small}
\begin{sc}
\begin{tabular}{lccc}
\toprule
Method & Calibration Samples & GPU Compute Time & Extra Optimization parameters \\
\midrule
Weight Averaging (WA) & 0 & $<$ 1 ms & None \\
Task Arithmetic (TA) & 0 & $<$ 1 ms & None \\
S-IPR \cite{paper5} & 0 & 14.2 seconds (SVD) & SVD projections \\
WCPR \cite{paper9} & 0 & 42.5 seconds (OT) & 1D OT transport plans \\
QR-IPR \cite{paper5} & 0 & 12.8 seconds & Parameter resonance scaling \\
\midrule
\textbf{DE-BN (Ours)} & \textbf{32} & \textbf{$<$ 2 ms} & \textbf{None} \\
\textbf{DE-QC (Ours)} & \textbf{32} & \textbf{1.8 seconds} & \textbf{Per-channel scaling factors} \\
\bottomrule
\end{tabular}
\end{sc}
\end{small}
\end{center}
\end{table}

While S-IPR, WCPR, and QR-IPR require significant offline GPU computation (up to 42 seconds for WCPR's optimal transport solvers) to align parameters, they fail to solve the core activation normalization pathology in low-bit regimes unless they are combined with a proper calibration baseline. In contrast, our proposed DE-BN updates statistics in under 2 milliseconds on a single GPU, and DE-QC vectorizes activation-MSE searches, completing the entire network's clipping scale calibration in under 2 seconds.

\subsection{Experimental Results under Robustness Environments (Noise and Blur)}
\label{app:robustness_results}
To supplement the discussion in Section 6.3, we provide the complete breakdown of multitask merging performance under Gaussian Noise and Defocus Blur across FP32, INT8, and INT4 precision regimes with and without BatchNorm calibration in Table~\ref{tab:robustness_results}.

Our findings under real-world corruptions mirror those of the clean environment:
\begin{enumerate}
    \item \textbf{Significant Noise and Blur Healing}: In the presence of Gaussian Noise, uncalibrated Tuned TA achieves only 39.70\% (FP32) and 39.28\% (INT4) accuracy. Our simple DE-BN with 32 samples restores accuracy to \textbf{77.22\%} (FP32) and \textbf{76.53\%} (INT4). Similarly, under Defocus Blur, DE-BN restores uncalibrated INT4 performance from 26.05\% to \textbf{55.60\%}.
    \item \textbf{Equivalence of Complex Algorithms}: In noisy environments, complex optimal transport alignment (WCPR) with 32 calibration samples achieves 75.63\% (INT4), while standard Tuned TA with DE-BN achieves \textbf{76.53\%}. This confirms that even under heavy input noise and low-bit quantization, proper normalization is the primary bottleneck, and complex weight-space alignment is redundant.
\end{enumerate}

\begin{table*}[htbp]
\caption{Average multitask accuracy on ResNet-18 across MNIST, Fashion-MNIST, and CIFAR-10 under Gaussian Noise and Defocus Blur environments for all combinations of precision regimes (FP32, INT8, INT4) and BN calibration sizes (0 vs 32 samples).}
\label{tab:robustness_results}
\begin{center}
\setlength{\tabcolsep}{4pt}
\begin{scriptsize}
\begin{sc}
\begin{tabular}{lcccccccccc}
\toprule
 & & & \multicolumn{4}{c}{\textbf{Gaussian Noise}} & \multicolumn{4}{c}{\textbf{Defocus Blur}} \\
\cmidrule(r){4-7} \cmidrule(l){8-11}
Method & Bits & BN-Cal & MNIST & FMNIST & CIFAR10 & Average & MNIST & FMNIST & CIFAR10 & Average \\
\midrule
\textbf{Precision: FP32} & & & & & & & & & & \\
WA & 32 & 0 & 38.62\% & 25.54\% & 14.90\% & 26.35\% & 24.63\% & 14.43\% & 16.88\% & 18.65\% \\
WA & 32 & 32 & 91.54\% & 63.85\% & 59.95\% & 71.78\% & 79.46\% & 59.24\% & 24.39\% & 54.36\% \\
Tuned TA & 32 & 0 & 64.78\% & 37.27\% & 17.04\% & 39.70\% & 38.79\% & 24.60\% & 17.74\% & 27.04\% \\
Tuned TA & 32 & 32 & 92.16\% & 80.90\% & 58.59\% & 77.22\% & 69.70\% & 60.64\% & 29.22\% & 53.19\% \\
Tuned TA + DE-QC & 32 & 0 & 64.78\% & 37.27\% & 17.04\% & 39.70\% & 38.79\% & 24.60\% & 17.74\% & 27.04\% \\
Tuned TA + DE-QC & 32 & 32 & 92.23\% & 76.42\% & 57.46\% & 75.37\% & 76.48\% & 59.55\% & 29.80\% & 55.28\% \\
S-IPR & 32 & 0 & 88.39\% & 50.77\% & 25.88\% & 55.01\% & 66.21\% & 28.13\% & 21.34\% & 38.56\% \\
S-IPR & 32 & 32 & 95.84\% & 75.50\% & 50.03\% & 73.79\% & 68.96\% & 59.29\% & 22.35\% & 50.20\% \\
WCPR & 32 & 0 & 91.79\% & 65.95\% & 29.14\% & 62.29\% & 70.79\% & 36.84\% & 23.46\% & 43.70\% \\
WCPR & 32 & 32 & 94.25\% & 67.81\% & 58.40\% & 73.49\% & 70.89\% & 63.93\% & 28.61\% & 54.48\% \\
QR-IPR & 32 & 0 & 87.43\% & 50.46\% & 24.89\% & 54.26\% & 58.29\% & 32.72\% & 20.65\% & 37.22\% \\
QR-IPR & 32 & 32 & 96.93\% & 77.86\% & 52.61\% & 75.80\% & 76.19\% & 51.78\% & 23.55\% & 50.51\% \\
\midrule
\textbf{Precision: INT8} & & & & & & & & & & \\
WA & 8 & 0 & 38.63\% & 25.22\% & 14.99\% & 26.28\% & 24.74\% & 14.33\% & 16.91\% & 18.66\% \\
WA & 8 & 32 & 89.45\% & 64.04\% & 59.04\% & 70.84\% & 75.75\% & 64.11\% & 30.32\% & 56.73\% \\
Tuned TA & 8 & 0 & 64.68\% & 37.23\% & 17.23\% & 39.71\% & 38.57\% & 24.70\% & 17.86\% & 27.04\% \\
Tuned TA & 8 & 32 & 95.37\% & 75.09\% & 57.53\% & 76.00\% & 78.49\% & 64.53\% & 31.18\% & 58.07\% \\
Tuned TA + DE-QC & 8 & 0 & 64.68\% & 37.23\% & 17.23\% & 39.71\% & 38.57\% & 24.70\% & 17.86\% & 27.04\% \\
Tuned TA + DE-QC & 8 & 32 & 91.72\% & 73.50\% & 57.70\% & 74.31\% & 75.67\% & 61.31\% & 29.34\% & 55.44\% \\
S-IPR & 8 & 0 & 88.54\% & 50.70\% & 25.79\% & 55.01\% & 66.03\% & 28.20\% & 21.27\% & 38.50\% \\
S-IPR & 8 & 32 & 92.66\% & 76.19\% & 49.34\% & 72.73\% & 80.09\% & 50.52\% & 26.95\% & 52.52\% \\
WCPR & 8 & 0 & 91.66\% & 65.96\% & 29.12\% & 62.25\% & 70.70\% & 36.53\% & 23.48\% & 43.57\% \\
WCPR & 8 & 32 & 96.86\% & 80.76\% & 55.53\% & 77.72\% & 79.55\% & 41.65\% & 28.57\% & 49.92\% \\
QR-IPR & 8 & 0 & 87.63\% & 50.59\% & 24.88\% & 54.37\% & 58.16\% & 32.94\% & 20.65\% & 37.25\% \\
QR-IPR & 8 & 32 & 94.05\% & 79.03\% & 54.06\% & 75.71\% & 82.01\% & 51.02\% & 24.50\% & 52.51\% \\
\midrule
\textbf{Precision: INT4} & & & & & & & & & & \\
WA & 4 & 0 & 37.57\% & 32.53\% & 14.64\% & 28.25\% & 19.52\% & 22.14\% & 16.90\% & 19.52\% \\
WA & 4 & 32 & 87.61\% & 79.03\% & 59.03\% & 75.22\% & 79.79\% & 60.34\% & 26.96\% & 55.70\% \\
Tuned TA & 4 & 0 & 61.28\% & 38.90\% & 17.65\% & 39.28\% & 34.26\% & 26.18\% & 17.72\% & 26.05\% \\
Tuned TA & 4 & 32 & 94.06\% & 75.47\% & 60.05\% & 76.53\% & 78.20\% & 62.91\% & 25.68\% & 55.60\% \\
Tuned TA + DE-QC & 4 & 0 & 61.28\% & 38.90\% & 17.65\% & 39.28\% & 34.26\% & 26.18\% & 17.72\% & 26.05\% \\
Tuned TA + DE-QC & 4 & 32 & 88.20\% & 76.22\% & 57.25\% & 73.89\% & 70.60\% & 60.67\% & 25.14\% & 52.14\% \\
S-IPR & 4 & 0 & 84.51\% & 47.77\% & 25.13\% & 52.47\% & 62.63\% & 25.29\% & 22.13\% & 36.68\% \\
S-IPR & 4 & 32 & 91.90\% & 71.11\% & 50.99\% & 71.33\% & 63.14\% & 45.23\% & 22.41\% & 43.59\% \\
WCPR & 4 & 0 & 89.18\% & 64.59\% & 31.24\% & 61.67\% & 67.29\% & 32.34\% & 25.78\% & 41.80\% \\
WCPR & 4 & 32 & 96.06\% & 77.07\% & 53.75\% & 75.63\% & 74.52\% & 64.13\% & 25.04\% & 54.56\% \\
QR-IPR & 4 & 0 & 85.84\% & 49.62\% & 24.74\% & 53.40\% & 55.14\% & 32.53\% & 21.73\% & 36.47\% \\
QR-IPR & 4 & 32 & 94.73\% & 71.04\% & 52.89\% & 72.89\% & 61.34\% & 60.62\% & 26.20\% & 49.39\% \\
\bottomrule
\end{tabular}
\end{sc}
\end{scriptsize}
\end{center}
\end{table*}

\subsection{Limitations and Practical Trade-offs}
\label{app:limitations}
As Methodologists, we believe that every scientific contribution must be accompanied by an honest, transparent deconstruction of its limitations, boundaries, and trade-offs. While our proposed DE-BN and DE-QC methods establish state-of-the-art baselines for quantized model merging and show that complex weight-alignment techniques are largely redundant, they possess specific constraints that dictate their practical application:

\begin{enumerate}
    \item \textbf{The Data-Efficient Assumption (Non-Zero Data):} 
    Our approach is explicitly ``data-efficient'' rather than ``data-free.'' It requires a small support set of 16 to 32 unlabeled samples per task to calibrate BatchNorm running statistics and weight clipping scales. 
    In extremely strict zero-resource deployment scenarios where absolutely zero unlabeled samples are available from the target task domain (e.g., due to severe regulatory or proprietary privacy restrictions), standard DE-BN and DE-QC cannot be executed. 
    However, we argue that the ``data-free'' constraint in prior model-merging works is often an artificial barrier that does not reflect real-world deployment. In virtually all practical edge scenarios, a device receives a stream of incoming inputs; accumulating a single batch of 32 unlabeled samples to perform a one-time calibration takes less than a second and requires no manual labeling. Thus, we view this data requirement as an extremely mild and highly practical trade-off for restoring full model accuracy and preventing catastrophic representation collapse.

    \item \textbf{Dependency on Task Identity (Task Routing):}
    Because DE-BN estimates running statistics specifically for each task, the model must be aware of the task context (the task identity $T$) during inference to load the corresponding task-specific BatchNorm statistics. 
    In mixed-input environments where samples from different tasks are fed to the model in an interleaved, sample-by-sample fashion without any task labels, loading task-specific statistics becomes challenging. 
    To deploy our method in such scenarios, a lightweight routing module or task-prediction head must be trained to dynamically select the appropriate BN parameters. Alternatively, an ensemble of task statistics can be combined, though this increases memory overhead. Exploring task-agnostic calibration that remains robust under 4-bit quantization is an exciting avenue for future work.

    \item \textbf{Calibration Overhead and Activation Caching:}
    For DE-QC, although the vectorized per-channel threshold optimization completes in under 2 seconds on GPU (Table~\ref{tab:efficiency_metrics}), it requires caching the FP32 activations of 32 samples across all layers. For very large models (such as modern LLMs), caching intermediate activations can lead to high peak memory usage during the brief calibration phase. While this is easily manageable on any modern edge device or cloud server, practitioners must budget the peak GPU/accelerator memory during calibration to prevent Out-Of-Memory (OOM) errors.
\end{enumerate}

In summary, while DE-BN and DE-QC do not satisfy the absolute zero-data constraint of pure zero-shot merging, they expose the massive performance-overhead trade-offs of existing data-free methods and establish that a tiny, practical amount of unlabeled data yields a massive increase in soundness, reliability, and accuracy under low-bit quantization.

\end{document}
